\documentclass[10pt,a4paper]{amsart}
\usepackage[margin=19mm]{geometry}
\usepackage{amsmath,amssymb,mathtools}
\usepackage[scr=rsfs,cal=euler]{mathalfa}
\usepackage{xfrac}
\usepackage[unicode,hidelinks,bookmarks=false]{hyperref}
\newcommand{\ie}{i.e.\ }
\newcommand{\cf}{cf.\ }
\newcommand{\resp}{respectively\ }
\newcommand{\Subsection}{\subsection}
\providecommand{\nobreakdash}{\nobreak\hbox{-}}
\setlength{\emergencystretch}{2em}
\multlinegap0pt
\usepackage{bm}
\usepackage{bbm}
\usepackage{dsfont}
\usepackage{xspace}
\usepackage{cancel}
\usepackage{mathtools}
\usepackage{amsthm}
\makeatletter
\@ifundefined{theorem}{\newtheorem{theorem}{Theorem}}{}
\@ifundefined{lemma}{\newtheorem{lemma}[theorem]{Lemma}}{}
\makeatother
\def \text{\mbox}
\def\C{\mathcal{C}}
\def\D{\mathcal{D}}
\def\E{\mathbb{E}}
\def\M{\mathcal{M}}
\def\X{\mathcal{X}}
\def\Y{\mathcal{Y}}
\def\Z{\mathcal {Z}}
\providecommand{\add}{\mathop{\rm add}\nolimits}%
\title[Mutation of $n$-cotorsion pairs in extriangulated categories]{Mutation of $n$-cotorsion pairs in extriangulated categories}
\author{Huimin Chang, Yu Liu, Panyue Zhou}
\date{Source: arXiv:2506.16188v1 (CC BY 4.0)}
\begin{document}
\maketitle

\noindent\emph{Source and licence.} Mutation of $n$-cotorsion pairs in extriangulated categories. Version arXiv:2506.16188v1 (CC BY 4.0), distributed under the Creative Commons Attribution 4.0 International licence, as recorded in the arXiv licence metadata for this exact version and shown on the abstract page https://arxiv.org/abs/2506.16188v1. Full rights notice: licenses/arxiv-2506.16188-CC-BY-4.0.txt.\par\smallskip

\noindent\textit{Editorial scope.} Counted unit: the Lemma whose source label is \texttt{a} in the article (source0.tex lines 803--810, source label \texttt{a}), delivered together with the article's own complete original proof of it (source0.tex lines 812--838, one \texttt{proof} environment of 27 lines). No mathematics is written, repaired or re-derived: statement and proof are the article's own text line for line. Required context retained uncounted: 4 (lines 782--787). Every definition, display and label of those lines is retained unchanged. Omitted from this package and not redistributed: 1--781, 788--802, 811--811, 839--1382. Nothing inside a retained excerpt is omitted or altered: every retained body is delivered exactly as the article's lines are written, so no omission receipt of any kind accompanies this package. Citations: the retained excerpts carry no citation command at all, so no bibliography entry of the article is transcribed and no reference is redistributed. Reference adaptation: none was needed and none was made; every reference inside the delivered unit resolves inside it, so no unresolved reference or citation is produced. Printed slips kept literal: no printed word, symbol, hypothesis or number of the counted statement or of its proof was altered. Layout and class: the article's own class is \texttt{amsart} with packages including amsmath, amssymb, amsthm, latexsym, xypic, amsfonts, mathrsfs, graphicx, color, xcolor, tikz, graphics, hyperref, anysize; the wrapper is the lane's pinned \texttt{amsart} editorial wrapper at 10pt on A4, which restates the theorem-like environments the retained text needs and adds the article's own macro definitions that the retained text needs (\texttt{\textbackslash C}, \texttt{\textbackslash D}, \texttt{\textbackslash E}, \texttt{\textbackslash M}, \texttt{\textbackslash X}, \texttt{\textbackslash Y}, \texttt{\textbackslash Z}, \texttt{\textbackslash add}, \texttt{\textbackslash text}), with extra wrapper packages (bm, bbm, dsfont, xspace, cancel, mathtools). The delivered numbering is the unit's own. Assets: no retained span contains a figure environment, a graphics-inclusion command, a PDF-page inclusion, a table or a TikZ picture; no image file is copied into this package, the package declares no asset dependency and no image file is redistributed with it. Licence: version arXiv:2506.16188v1 of this article is distributed under the Creative Commons Attribution 4.0 International licence, as recorded in the arXiv licence metadata for this exact version and shown on the abstract page https://arxiv.org/abs/2506.16188v1. A full rights notice is delivered at licenses/arxiv-2506.16188-CC-BY-4.0.txt. Characters: every retained excerpt is pure ASCII, so the delivered engine, XeLaTeX with its default Latin Modern text font, reports no missing character.
\begin{lemma}\label{4}
Let $\M_i$ be a contravariantly finite and extension closed subcategory of an extriangulated category $\C$ such that
$\E(\M_i,\M_j)=0$ for any $i<j$. Put
$$\X_n\colon=\add(\M_1\ast\M_2\ast\cdots\ast\M_n),\;\;\Y_n=\colon\bigcap\limits_{k=1}^{n}\M_i^\perp.$$
Then $(\X_n,\Y_n)$ is a torsion pair.
\end{lemma}

\begin{lemma}\label{a}
Let $\X$ be a subcategory of $\C$ satisfying  $\D\subset\X\subset\Z$. Then
\begin{itemize}
  \item [\rm (1)] $\X$ is a contravariantly finite subcategory of $\C$ if and only if $\overline{\X}$ is a contravariantly finite subcategory of $\mathfrak{U}$.
  \item [\rm (2)] $\X$ is a covariantly finite subcategory of $\C$ if and only if $\overline{\X}$ is a covariantly finite subcategory of $\mathfrak{U}$.
  \item [\rm (3)] $\X$ is a functorially finite subcategory of $\C$ if and only if $\overline{\X}$ is a functorially finite subcategory of $\mathfrak{U}$.
\end{itemize}
\end{lemma}

\begin{proof}
We only prove the first assertion, the others can be proved similarly.

Assume that $\X$ is a contravariantly finite subcategory of $\C$. Let $Z$ be an object in $\Z$. Since $\X$ is contravariantly finite, there exists an object $X \in \X$ and a morphism $f \colon X \to Z$ which serves as a right $\X$-approximation of $Z$ in $\C$. It is straightforward to verify that the induced morphism $\overline{f} \colon X \to Z$ is a right $\overline{\X}$-approximation of $Z$ in the quotient category $\mathfrak{U}$.

Conversely, assume that $\overline{\X}$ is a contravariantly finite subcategory of $\mathfrak{U}$. Since $\D$ is functorially finite in $\C$, it follows from Lemma~\ref{4} and its dual that $\Z$ is also functorially finite in $\C$. Let $C$ be an object in $\C$. Then, as $\Z$ is functorially finite, there exists an object $Z \in \Z$ along with a morphism $f \colon Z \to C$ that is a right $\Z$-approximation of $C$.

By assumption, there exists $X \in \X$ and a morphism $\overline{g} \colon X \to Z$ in $\mathfrak{U}$ such that it is a right $\overline{\X}$-approximation of $Z$. Moreover, since $\D$ is functorially finite in $\C$ and $Z \in \Z \subseteq \C$, there exists an object $D \in \D$ with a morphism $c \colon D \to Z$ forming a right $\D$-approximation of $Z$. We claim that the morphism
\[
(f \circ g,\, f \circ c) \colon X \oplus D \to C
\]
is a right $\X$-approximation of $C$ in $\C$.

To verify the claim, consider any morphism $h \colon X' \to C$ with $X' \in \X \subseteq \Z$. Since $f \colon Z \to C$ is a right $\Z$-approximation, there exists $\ell \colon X' \to Z$ such that $h = f \circ \ell$. As $\overline{g} \colon X \to Z$ is a right $\overline{\X}$-approximation in $\mathfrak{U}$, there exists a morphism $\overline{k} \colon X' \to X$ such that $\overline{\ell} = \overline{g} \circ \overline{k}$ and then
$
\overline{\ell - g \circ k} = \overline{0}.
$
Hence, the morphism $\ell - g \circ k$ factors through some $D_1 \in \D$; that is, there exist morphisms $a_1 \colon X' \to D_1$ and $a_2 \colon D_1 \to Z$ such that $\ell - g \circ k = a_2 \circ a_1$. Since $c \colon D \to Z$ is a right $\D$-approximation, there exists $d \colon D_1 \to D$ such that $a_2 = c \circ d$. Consequently,
\[
\ell = g \circ k + c \circ d \circ a_1 \quad \text{and} \quad h = f \circ \ell = f \circ (g \circ k + c \circ d \circ a_1).
\]
This shows that $h$ factors through the morphism $(f \circ g, f \circ c) \colon X \oplus D \to C$ via
\[
X' \xrightarrow{\binom{k}{d \circ a_1}} X \oplus D.
\]
Thus, we conclude that $(f \circ g, f \circ c)$ is indeed a right $\X$-approximation of $C$ in $\C$, and therefore, $\X$ is contravariantly finite.
\end{proof}

\end{document}
